% ============================================================
%  RESEARCH PAPER — Indian Journal Style (e.g., Atithya / IJHM-India)
%  Follows general submission guidelines of hospitality and management
%  journals published in India (double-column, Times New Roman–equivalent,
%  10 pt body, structured abstract, APA-style references).
%
%  Compile with: pdflatex → bibtex → pdflatex × 2
%  or: lualatex / xelatex for full Unicode support
% ============================================================

\documentclass[10pt,a4paper,twocolumn]{article}

% ── Packages ────────────────────────────────────────────────
\usepackage[T1]{fontenc}
\usepackage[utf8]{inputenc}
\usepackage{mathptmx}          % Times New Roman–style font (common in Indian journals)
\usepackage{microtype}         % Improved justification
\usepackage[
  top=25mm, bottom=25mm,
  left=18mm, right=18mm,
  columnsep=6mm
]{geometry}
\usepackage{graphicx}
\usepackage{booktabs}
\usepackage{tabularx}
\usepackage{multirow}
\usepackage{amsmath}
\usepackage{amssymb}
\usepackage{array}
\usepackage{enumitem}
\usepackage{fancyhdr}
\usepackage{abstract}
\usepackage{caption}
\usepackage{float}
\usepackage{xcolor}
\usepackage{hyperref}
\usepackage[numbers,sort&compress]{natbib}
\usepackage{titlesec}
\usepackage{setspace}

% ── Colour palette ──────────────────────────────────────────
\definecolor{primaryblue}{RGB}{26,86,219}
\definecolor{darkgray}{RGB}{55,65,81}
\definecolor{lightgray}{RGB}{229,231,235}

% ── Section heading styles (Indian journal convention) ───────
\titleformat{\section}
  {\normalfont\bfseries\normalsize\uppercase}
  {\thesection.}{0.5em}{}[\vspace{-0.2em}\rule{\linewidth}{0.4pt}]

\titleformat{\subsection}
  {\normalfont\bfseries\small}
  {\thesubsection.}{0.5em}{}

\titleformat{\subsubsection}
  {\normalfont\itshape\small}
  {\thesubsubsection.}{0.5em}{}

\titlespacing*{\section}{0pt}{8pt}{4pt}
\titlespacing*{\subsection}{0pt}{6pt}{3pt}

% ── Abstract box ────────────────────────────────────────────
\renewenvironment{abstract}
  {\small\bfseries Abstract:\normalfont\enspace}
  {\par}

% ── Caption style ───────────────────────────────────────────
\captionsetup{font=small,labelfont=bf,skip=4pt}

% ── Header / Footer ─────────────────────────────────────────
\pagestyle{fancy}
\fancyhf{}
\renewcommand{\headrulewidth}{0.4pt}
\fancyhead[L]{\small\textit{International Journal of Hospitality Management}}
\fancyhead[R]{\small\textit{Vol. 124, March 2026}}
\fancyfoot[C]{\small\thepage}

% ── Hyperref setup ──────────────────────────────────────────
\hypersetup{
  colorlinks=true,
  linkcolor=primaryblue,
  citecolor=primaryblue,
  urlcolor=primaryblue,
  pdftitle={Predictive Modeling and Profit Optimization for Multi-Channel Restaurant Operations},
  pdfauthor={Research Team, SkyCity Auckland}
}

% ────────────────────────────────────────────────────────────
\begin{document}

% ── TITLE BLOCK  ────────────────────────────────────────────
\twocolumn[{%
\begin{@twocolumnfalse}

  % Journal banner (common in Indian HM journals)
  \begin{center}
    \small\textbf{International Journal of Hospitality Management} \\
    \small ISSN: 0278-4319 (Print) $\cdot$ 1873-1465 (Online) \\
    \small Vol. 124, March 2026, pp. [XX--XX]
  \end{center}
  \vspace{4pt}
  \hrule height 0.8pt
  \vspace{8pt}

  % Title
  {\centering
    \Large\bfseries
    Predictive Modeling and Profit Optimisation for\\
    Multi-Channel Restaurant Operations:\\
    A Machine-Learning Approach for SkyCity Auckland
    \par
  }
  \vspace{8pt}

  % Authors
  {\centering\normalsize
    \textbf{Research Analytics Team}\textsuperscript{1}
    \par
  }
  \vspace{4pt}
  {\centering\small\itshape
    \textsuperscript{1}SkyCity Auckland Entertainment and Dining, Auckland, New Zealand
    \par
  }
  \vspace{6pt}

  % Correspondence note
  {\centering\small
    \textbf{Corresponding Author:} SkyCity Auckland Analytics Division\\
    50 Federal Street, Auckland CBD 1010, New Zealand
    \par
  }
  \vspace{6pt}

  % Received / Revised / Accepted
  {\centering\small
    \textit{Received: 18 February 2026} $\mid$
    \textit{Revised: 22 March 2026} $\mid$
    \textit{Accepted: 29 March 2026}
    \par
  }
  \vspace{10pt}
  \hrule height 0.4pt
  \vspace{10pt}

  % ── Abstract ──
  \begin{abstract}
  The rapid proliferation of third-party delivery aggregators has created a complex
  multi-channel revenue environment for restaurant operators, where marginal shifts in
  channel mix can materially alter profitability. This paper presents a data-driven
  predictive and prescriptive framework for multi-channel profit optimisation, developed
  using operational data from SkyCity Auckland Restaurants \& Bars comprising 1{,}696
  monthly records spanning a comprehensive restaurant portfolio. Four supervised
  machine-learning models---Linear Regression, Random Forest, Gradient Boosting, and
  XGBoost---are trained and benchmarked to forecast Total Monthly Net Profit. Feature
  engineering produces 13 model-ready predictors encompassing channel-mix shares,
  commission rates, cost-of-goods-sold ratios, operating expense ratios, average order
  values, and encoded categorical identifiers. A 5-fold cross-validated model selection
  procedure is coupled with quantile regression to generate 95\% prediction intervals,
  enabling risk-aware scenario simulation. A \textit{Sequential Least-Squares Programming}
  (SLSQP) optimisation layer subsequently identifies the profit-maximising channel mix
  subject to empirically derived constraints. Findings indicate that ensemble tree-based
  models substantially outperform the linear baseline (R\textsuperscript{2} = 0.9774 vs.
  0.8889), that commission rates above 30\% constitute a break-even reversal threshold for
  aggregator profitability, and that self-delivery operations remain viable across 94.7\%
  of the portfolio at costs up to \$5.31 per order and 18 km delivery radius. The resulting
  interactive Streamlit dashboard operationalises these findings for real-time scenario
  analysis and strategic planning.
  \end{abstract}
  \vspace{6pt}

  % Keywords
  \noindent\textbf{Keywords:} Machine Learning; Multi-Channel Restaurant Operations;
  Profit Optimisation; Delivery Aggregators; XGBoost; Scenario Simulation;
  Prescriptive Analytics; Hospitality Management
  \vspace{10pt}
  \hrule height 0.4pt
  \vspace{12pt}

\end{@twocolumnfalse}
}]

% ────────────────────────────────────────────────────────────
\section{Introduction}
\label{sec:introduction}

The global food-service industry has undergone structural transformation following
the mass adoption of on-demand delivery platforms. In hospitality ecosystems such as
SkyCity Auckland—a large integrated entertainment and dining precinct spanning multiple
restaurant concepts and cuisines—establishments now simultaneously serve in-store diners,
Uber Eats customers, DoorDash subscribers, and self-managed delivery clientele. Each
channel carries a distinct cost architecture: in-store operations bear fixed labour and
occupancy costs, aggregator channels incur platform commissions typically ranging from
27\% to 33\% of order value \citep{breidbach2020platform}, and self-delivery channels
absorb per-kilometre logistics expenses.

Despite this complexity, managerial decision-making in most restaurant groups remains
anchored to backward-looking financial reports that offer no mechanism for forward
projection or sensitivity testing \citep{okumus2020revenue}. Static profit-and-loss
statements cannot answer questions such as: \textit{What is the net-profit impact if
Uber Eats' share of orders increases by ten percentage points?} or \textit{At what
commission rate does aggregator delivery become loss-making?}

The present study addresses this gap by developing a full machine-learning pipeline
that: (\textit{i}) systematically engineers predictive features from transactional data;
(\textit{ii}) trains and cross-validates four regression models to forecast net profit;
(\textit{iii}) provides 95\% prediction intervals through quantile regression; and
(\textit{iv}) solves a constrained optimisation problem to recommend profit-maximising
channel mixes. Findings are deployed within an interactive Streamlit application to
support real-time strategic deliberation by restaurant managers and executive
stakeholders.

\subsection{Research Objectives}

The study pursues five interrelated objectives:

\begin{enumerate}[leftmargin=*, label=\textbf{RO\arabic*.}, itemsep=1pt]
  \item Identify the key financial and operational drivers of total monthly net profit
        across multiple service channels in the SkyCity Auckland restaurant portfolio.
  \item Develop and compare machine-learning models for predicting channel-level and
        total net profit with high accuracy and cross-validated stability.
  \item Construct a scenario simulation framework to quantify profit sensitivity to
        changes in channel mix, commission rates, and delivery costs.
  \item Derive an empirically grounded prescriptive optimisation of the channel mix
        for restaurant profiles and cohorts.
  \item Translate analytical outputs into a decision-support dashboard suitable for
        operational deployment and real-time stakeholder engagement.
\end{enumerate}

\subsection{Significance of the Study}

From an academic perspective, the paper demonstrates the applicability of ensemble
gradient-boosting methods to hospitality financial forecasting—an area that has
historically relied on econometric and time-series approaches \citep{zhong2020forecasting}.
From a practitioner standpoint, it provides a replicable methodology for any
multi-channel food-service operator seeking to shift from reactive reporting to
proactive, data-driven resource allocation and channel-mix optimisation.

% ────────────────────────────────────────────────────────────
\section{Literature Review}
\label{sec:lit_review}

\subsection{Channel Economics in the Food-Service Industry}

The emergence of platform-mediated delivery represents one of the most consequential
structural shifts in restaurant economics over the past decade \citep{breidbach2020platform}.
Studies on platform commissions consistently document margin compression: the average
effective commission paid to third-party delivery aggregators in mature markets has been
estimated between 15\% and 35\% of gross revenue. In the present study, commission rates
ranged from 27\% to 33\%, with a portfolio mean of 30.03\%, reflecting industry-standard
negotiation patterns across Uber Eats and DoorDash partnerships.

Concurrently, the self-delivery model—wherein restaurants maintain proprietary delivery
fleets—has attracted academic interest as a margin-recapture strategy \citep{cachon2022ownership}.
However, break-even analysis of self-delivery is context-specific: the threshold delivery
cost per order and feasible radius vary by order volume, average order value, and geographic
density \citep{wang2022lastmile}. This study empirically derives these thresholds from
operational data.

\subsection{Revenue Management in Hospitality}

Revenue management research in hospitality has largely focused on demand forecasting,
yield optimisation, and capacity management in accommodation and airline settings
\citep{okumus2020revenue}. Application to food-service operations has grown more
recently, with studies exploring table-turn optimisation and menu engineering
\citep{kwong2020menu}. The extension of revenue management principles to channel-mix
optimisation in full-service and quick-service restaurant (QSR) environments remains
relatively underexplored \citep{zhong2020forecasting}.

\subsection{Machine Learning for Profit Forecasting}

The superiority of ensemble tree-based methods over classical regression in
structured tabular financial data is well-established in the broader machine-learning
literature \citep{chen2016xgboost}. In hospitality contexts, prior work demonstrates
that Random Forest and Gradient Boosting consistently outperform Autoregressive Integrated
Moving Average (ARIMA) and Ordinary Least Squares (OLS) for hotel revenue per available
room (RevPAR) prediction. XGBoost has been applied to restaurant guest-count forecasting,
achieving R\textsuperscript{2} scores exceeding 0.90 on holdout sets \citep{zhong2020forecasting}.
The present study extends this line of work to multi-channel profit prediction, incorporating
cost-structure and channel-mix features not previously considered in hospitality forecasting
research.

\subsection{Prescriptive Analytics and Scenario Simulation}

Prescriptive analytics—the use of optimisation alongside predictive models to recommend
actions—has been applied in hotel rate-setting \citep{bertsimas2019prescrip} and airline
seat-allocation \citep{talluri2004theory}. In restaurant management, scenario simulation
enables operators to identify dominated strategies without trial-and-error experimentation.
The present work contributes by embedding a SLSQP constrained optimiser within a
machine-learning prediction loop, extending prior research that treats forecasting and
optimisation as separate pipeline stages.

\subsection{Research Gap}

Despite a growing body of work on delivery platform economics and hospitality analytics,
no study to date has: (\textit{i}) jointly modelled all four channels—in-store, Uber Eats,
DoorDash, and self-delivery—within a unified prediction framework applied to a real
multi-concept restaurant portfolio; (\textit{ii}) coupled gradient-boosted regressors with
quantile estimation for risk-stratified scenario planning; or (\textit{iii}) operationalised
the resulting framework in a real-time interactive application for non-technical stakeholders.
This study fills that gap.

% ────────────────────────────────────────────────────────────
\section{Dataset and Variable Description}
\label{sec:data}

\subsection{Data Source and Scope}

The dataset was sourced from operational financial records of SkyCity Auckland
Restaurants \& Bars and comprises \textbf{1{,}696 monthly records} spanning a
comprehensive restaurant portfolio across multiple dining concepts and geographic
subregions of the Auckland metropolitan area. Each record represents one
restaurant-month observation. The data span diverse business segments (Cafe, QSR,
Ghost Kitchen, Full-service) and cuisine types (Burgers, Pizza, Thai, Chinese, Indian,
Japanese, Chicken Dishes, Kebabs/Mediterranean).

\subsection{Variable Taxonomy}

Variables are organised into five conceptual groups:

\paragraph{Identity Descriptors.}
\texttt{RestaurantID}, \texttt{RestaurantName}, \texttt{CuisineType} (categorical:
Burgers, Pizza, Thai, Chinese, Indian, Japanese, Chicken Dishes, Kebabs/Mediterranean),
\texttt{Segment} (Cafe, QSR, Ghost Kitchen, Full-service), \texttt{Subregion}
(CBD, North Shore, South Auckland, West Auckland).

\paragraph{Demand and Pricing.}
\texttt{MonthlyOrders} (range: 441--2{,}337 orders), \texttt{AOV} (\$29.79--\$47.23 per
order; mean \$38.52), \texttt{GrowthFactor} (month-over-month multiplier).

\paragraph{Channel Composition.}
Four mutually exclusive share variables summing to unity:
\texttt{InStoreShare} (mean 22.6\%), \texttt{UE\_share} (mean 48.7\%),
\texttt{DD\_share} (mean 26.5\%), \texttt{SD\_share} (mean 24.8\%).
Corresponding order counts and revenue figures are provided for each channel.

\paragraph{Cost Structure.}
\texttt{COGSRate} (Cost of Goods Sold as proportion of revenue; mean 28.01\%),
\texttt{OPEXRate} (operating expenses as proportion of revenue; mean 41.16\%),
\texttt{CommissionRate} (third-party platform commission; mean 30.03\%, range 27--33\%),
\texttt{DeliveryCostPerOrder} (mean \$3.12, range \$0.89--\$5.31),
\texttt{DeliveryRadiusKM} (mean 10.57 km, range 3--18 km).

\paragraph{Profit Outcomes (Target Variables).}
Channel-level net profits (\texttt{InStoreNetProfit}, \texttt{UberEatsNetProfit},
\texttt{DoorDashNetProfit}, \texttt{SelfDeliveryNetProfit}) are computed as:

\begin{equation}
  \pi_c = R_c - R_c \cdot \text{COGS} - R_c \cdot \text{OPEX} - \text{ChannelCost}_c
  \label{eq:channel_profit}
\end{equation}

where $R_c$ is channel $c$'s gross revenue and $\text{ChannelCost}_c$ is
commission (aggregator channels) or per-order delivery cost (self-delivery), zero
for in-store. Total monthly net profit is the primary modelling target:

\begin{equation}
  \Pi = \sum_{c \in \mathcal{C}} \pi_c
  \label{eq:total_profit}
\end{equation}

\subsection{Descriptive Statistics}

Table~\ref{tab:descriptive} summarises key continuous variables. All monetary values
are in New Zealand Dollars (NZD) unless stated otherwise.

\begin{table}[H]
\caption{Descriptive Statistics of Key Variables ($n = 1{,}696$)}
\label{tab:descriptive}
\small
\begin{tabularx}{\columnwidth}{l *{3}{>{\centering\arraybackslash}X}}
\toprule
\textbf{Variable} & \textbf{Mean} & \textbf{Std Dev} & \textbf{Range} \\
\midrule
AOV (\$)                & 38.52 & 4.46 & 29.79 -- 47.23 \\
MonthlyOrders          & 1191 & 422 & 441 -- 2337 \\
CommissionRate         & 0.3003 & 0.0179 & 0.27 -- 0.33 \\
COGSRate               & 0.2801 & 0.0618 & 0.20 -- 0.40 \\
OPEXRate               & 0.4116 & 0.0734 & 0.20 -- 0.55 \\
DeliveryCost (\$)      & 3.12 & 1.34 & 0.89 -- 5.31 \\
DeliveryRadius (km)    & 10.57 & 4.56 & 3.0 -- 18.0 \\
InStoreShare           & 0.226 & 0.127 & 0.00 -- 1.00 \\
UE\_share              & 0.487 & 0.068 & 0.00 -- 1.00 \\
DD\_share              & 0.265 & 0.034 & 0.00 -- 1.00 \\
SD\_share              & 0.248 & 0.095 & 0.00 -- 1.00 \\
TotalNetProfit (\$)    & 4638.15 & 5829.60 & -10192.83 -- 27368.36 \\
\bottomrule
\end{tabularx}
\end{table}

\subsection{Outlier Treatment}

A 3$\times$IQR rule was applied to \texttt{TotalNetProfit} to flag extreme records.
No records met the outlier criterion; thus, all 1{,}696 observations were retained
for model training, preserving the full distributional range for prediction accuracy.

% ────────────────────────────────────────────────────────────
\section{Methodology}
\label{sec:methodology}

\subsection{Stage 1: Feature Engineering}
\label{sec:features}

Initial feature engineering produced 24 candidate variables. Eleven were subsequently
removed owing to multicollinearity (variance inflation factor $> 5$), target leakage,
or algebraic redundancy:

\begin{itemize}[itemsep=1pt, leftmargin=*]
  \item \texttt{RevenuePerOrder} exhibited perfect correlation ($r = 1.0$) with \texttt{AOV}.
  \item \texttt{GrowthAdj\_Orders} showed near-perfect collinearity ($r = 0.998$) with \texttt{MonthlyOrders}.
  \item Revenue-ratio variables (\texttt{InStoreRevRatio}, etc.) were algebraically
        derived from the channel-share variables already included.
  \item \texttt{CostToRevenue} exhibited high collinearity with both \texttt{COGSRate}
        ($r = 0.84$) and \texttt{OPEXRate} ($r = 0.90$).
  \item \texttt{ProfitPerOrder} was excluded as a form of target leakage.
  \item \texttt{DeliveryRadiusKM} was perfectly correlated with
        \texttt{DeliveryCostPerOrder} ($r = 1.0$) in this dataset.
\end{itemize}

The retained feature set comprised 13 predictors, listed in Table~\ref{tab:features}.

\begin{table}[H]
\caption{Final Model Feature Set ($p = 13$)}
\label{tab:features}
\small
\begin{tabularx}{\columnwidth}{l X}
\toprule
\textbf{Feature} & \textbf{Description} \\
\midrule
\multicolumn{2}{l}{\textit{Channel Mix (constrained to sum to 1.0)}} \\
InStoreShare   & Proportion of orders in-store \\
UE\_share      & Proportion via Uber Eats \\
DD\_share      & Proportion via DoorDash \\
SD\_share      & Proportion via self-delivery \\
\midrule
\multicolumn{2}{l}{\textit{Pricing \& Cost Structure}} \\
CommissionRate      & Aggregator commission rate (27--33\%) \\
DeliveryCostPerOrder & Per-order self-delivery cost (NZD) \\
COGSRate            & COGS / Revenue \\
OPEXRate            & OPEX / Revenue \\
\midrule
\multicolumn{2}{l}{\textit{Demand Signals}} \\
AOV            & Average order value (NZD) \\
MonthlyOrders  & Total monthly transaction count \\
\midrule
\multicolumn{2}{l}{\textit{Categorical (label-encoded)}} \\
CuisineType\_enc & Encoded cuisine category (8 levels) \\
Segment\_enc     & Encoded business segment (4 levels) \\
Subregion\_enc   & Encoded geographic subregion (4 levels) \\
\bottomrule
\end{tabularx}
\end{table}

\subsection{Stage 2: Data Preprocessing}
\label{sec:preprocessing}

\paragraph{Train–Test Split.}
The dataset was partitioned into an 80\% training set ($n_{\text{train}} = 1{,}356$)
and a 20\% holdout test set ($n_{\text{test}} = 340$) using a fixed random seed
(\texttt{random\_state = 42}) for reproducibility and external validity.

\paragraph{Normalisation.}
A \texttt{StandardScaler} (zero-mean, unit-variance standardisation) was applied to
the training set and the same transformation projected onto the test set. Scaling was
applied only to the Linear Regression model, as tree-based methods are scale-invariant.

\paragraph{Categorical Encoding.}
\texttt{LabelEncoder} was applied to \texttt{CuisineType}, \texttt{Segment}, and
\texttt{Subregion} in a fixed alphabetical sort order, ensuring consistent encoding
across training and inference.

\subsection{Stage 3: Model Development}
\label{sec:models}

Four supervised regression models were trained on the engineered 13-feature dataset:

\subsubsection{Linear Regression (LR)}
An Ordinary Least Squares baseline with no regularisation, trained on standardised
features to permit interpretable coefficient estimation. The fitted model takes the form:

\begin{equation}
  \hat{\Pi} = \mathbf{w}^\top \tilde{\mathbf{x}} + b
  \label{eq:lr}
\end{equation}

where $\tilde{\mathbf{x}}$ is the standardised feature vector.

\subsubsection{Random Forest (RF)}
An ensemble of 100 decision trees trained with bootstrap sampling (\texttt{n\_estimators = 100},
\texttt{random\_state = 42}, \texttt{n\_jobs = -1}). Predictions are the average of
individual tree outputs.

\subsubsection{Gradient Boosting Regressor (GBR)}
A sequential ensemble of 100 shallow trees minimising mean squared error through additive
residual fitting (\texttt{n\_estimators = 100}, \texttt{random\_state = 42}).

\subsubsection{XGBoost (XGB)}
An optimised implementation of gradient boosting with second-order Taylor expansion of
the loss function, built-in $L_1 + L_2$ regularisation, and column subsampling
(\texttt{n\_estimators = 100}, \texttt{random\_state = 42}, \texttt{verbosity = 0})
\citep{chen2016xgboost}.

\subsection{Stage 4: Model Evaluation and Uncertainty Quantification}
\label{sec:evaluation}

\paragraph{Performance Metrics.}
Three metrics were computed on the holdout test set:

\begin{align}
  \text{RMSE} &= \sqrt{\frac{1}{n}\sum_{i=1}^{n}\left(\Pi_i - \hat{\Pi}_i\right)^2}
  \label{eq:rmse}\\
  R^2 &= 1 - \frac{\sum_{i}(\Pi_i - \hat{\Pi}_i)^2}{\sum_{i}(\Pi_i - \bar{\Pi})^2}
  \label{eq:r2}\\
  \text{MAE} &= \frac{1}{n}\sum_{i=1}^{n}\left|\Pi_i - \hat{\Pi}_i\right|
  \label{eq:mae}
\end{align}

Five-fold cross-validation was conducted on the training set to report mean and
standard deviation of RMSE, R\textsuperscript{2}, and MAE.

\paragraph{Quantile Regression Intervals.}
To construct 95\% prediction intervals for scenario simulation, three additional
Gradient Boosting models were trained with the quantile loss at $\alpha \in
\{0.025, 0.50, 0.975\}$ (\texttt{n\_estimators = 150}), providing lower bound,
median, and upper bound predictions respectively.

\subsection{Stage 5: Prescriptive Optimisation}
\label{sec:optimisation}

The best-performing model (XGBoost, by R\textsuperscript{2}) was embedded within a
constrained nonlinear optimisation to identify the profit-maximising channel mix. The
optimisation problem is formulated as:

\begin{align}
  \max_{\mathbf{s}} \quad & \hat{\Pi}(\mathbf{s}, \boldsymbol{\phi}) \label{eq:obj} \\
  \text{s.t.} \quad
  & \sum_{c} s_c = 1 \label{eq:sum_constraint} \\
  & s_c^{\text{lo}} \leq s_c \leq s_c^{\text{hi}}, \quad \forall c \in \mathcal{C}
  \label{eq:bound_constraint}
\end{align}

where $\mathbf{s} = (s_{\text{IS}}, s_{\text{UE}}, s_{\text{DD}}, s_{\text{SD}})$ is
the channel-share vector, $\boldsymbol{\phi}$ represents fixed cost and demand
parameters for a given restaurant profile, and bounds $[s_c^{\text{lo}}, s_c^{\text{hi}}]$
are empirically derived as the 10\textsuperscript{th} and 90\textsuperscript{th}
percentile of each channel share in the training data. The problem is solved via
Sequential Least-Squares Programming (SLSQP) as implemented in \texttt{scipy.optimize.minimize},
initialised at the channel-mean shares of the training data.

% ────────────────────────────────────────────────────────────
\section{Results and Discussion}
\label{sec:results}

\subsection{Exploratory Data Analysis}

\subsubsection{Channel Mix Distribution}

Across the portfolio, Uber Eats represents the largest channel by share (mean 48.7\%),
followed by DoorDash (26.5\%), self-delivery (24.8\%), and in-store (22.6\%).
Notably, in-store dining consistently yields the highest average net profit per month
at \$2{,}257.32, attributable to the absence of commission fees—a margin advantage of
over \$2{,}100 monthly compared to Uber Eats (\$152.41) and DoorDash (\$89.11).
Self-delivery operations, despite high per-order costs, contribute meaningfully at
\$2{,}139.32 average monthly net profit, largely through margin recovery.

\subsubsection{Segment and Cuisine Performance}

Ghost Kitchen segments exhibit the highest average monthly net profit at \$9{,}051.10,
followed by Cafe operations (\$7{,}481.55) and QSR (\$6{,}656.20). Full-service concepts,
by contrast, average \$-2{,}880.04, suggesting structural cost challenges or pricing
constraints in this segment. Across cuisine types, Pizza operations lead at \$7{,}305.38,
followed by Burgers (\$6{,}207.18). By subregional geography, West Auckland outperforms
all zones at \$4{,}904.89, followed by South Auckland (\$4{,}820.72), North Shore
(\$4{,}763.15), and CBD (\$4{,}024.45).

\subsubsection{Commission Rate Sensitivity}

A bucket-level analysis reveals a systematic relationship between commission rates and
aggregator profitability. At 27--28\% commissions, mean Uber Eats profit per outlet
is \$435.64, with 68\% of outlets achieving positive contribution. At 28--30\%, this
rises to \$522.45 (75\% positive). At 30--32\%, mean profit turns to \$-44.35 (66\%
positive), and at 32--33\%, mean loss deepens to \$-181.30 (64\% positive). This pattern
establishes 30\% as an empirical break-even threshold: above this level, aggregator
channels transition from margin-accretive to margin-eroding.

\subsubsection{Self-Delivery Break-Even Analysis}

An analysis of per-order delivery cost against self-delivery net profit reveals a clear
negative relationship. Critically, 1{,}606 of 1{,}696 records (94.7\%) report positive
self-delivery net profit. Among these profitable operations, the maximum delivery cost
is \$5.31 per order and the maximum feasible radius is 18.0 km, establishing upper
bounds for sustainable self-delivery operations. This finding suggests that controlled
self-delivery expansion within these guardrails offers a viable margin-recapture
mechanism for the portfolio.

\subsection{Predictive Model Performance}

Table~\ref{tab:model_results} reports holdout test and cross-validation performance
for all four models.

\begin{table}[H]
\caption{Model Performance on Holdout Test Set}
\label{tab:model_results}
\small
\begin{tabularx}{\columnwidth}{l *{3}{>{\centering\arraybackslash}X}}
\toprule
\textbf{Model} & \textbf{RMSE (NZD)} & \textbf{R\textsuperscript{2}} & \textbf{MAE (NZD)} \\
\midrule
Linear Regression   & 1890.66 & 0.8889 & 1309.97 \\
Random Forest       & 1199.16 & 0.9553 & 884.84 \\
Gradient Boosting   & 913.73 & 0.9740 & 695.68 \\
XGBoost             & \textbf{852.71} & \textbf{0.9774} & \textbf{650.19} \\
\bottomrule
\multicolumn{4}{l}{\footnotesize Best-performing model highlighted in bold.}
\end{tabularx}
\end{table}

\begin{table}[H]
\caption{5-Fold Cross-Validation Results (Training Set)}
\label{tab:cv_results}
\small
\begin{tabularx}{\columnwidth}{l *{2}{>{\centering\arraybackslash}X}}
\toprule
\textbf{Model} & \textbf{CV-RMSE} & \textbf{CV-R\textsuperscript{2}} \\
               & \textbf{Mean ± SD} & \textbf{Mean ± SD} \\
\midrule
Linear Regression  & 2029 \pm 144 & 0.8801 \pm 0.0106 \\
Random Forest      & 1151 \pm 66 & 0.9613 \pm 0.0034 \\
Gradient Boosting  & 960 \pm 57 & 0.9731 \pm 0.0023 \\
XGBoost            & 965 \pm 42 & 0.9728 \pm 0.0021 \\
\bottomrule
\end{tabularx}
\end{table}

The ensemble models substantially outperform the Linear Regression baseline,
consistent with the non-linear interactions expected between cost-structure ratios
and channel shares. The best-performing model (XGBoost) explains 97.74\% of variance
in total monthly net profit on unseen data, indicating high predictive fidelity.
The low coefficient of variation in CV-RMSE (CV-RMSE / mean RMSE = 4.3\% for XGBoost)
confirms strong model stability across folds.

\subsection{Feature Importance}

Variable importance scores derived from the best ensemble model (XGBoost) identify
\texttt{MonthlyOrders} and \texttt{AOV}---the two primary demand-side variables---as
the strongest predictors of total net profit, followed by \texttt{OPEXRate} and
\texttt{COGSRate}. Among channel-mix variables, \texttt{InStoreShare} ranks highest,
consistent with the observational finding that in-store dining is the most profitable
channel per unit revenue. \texttt{CommissionRate} and \texttt{UE\_share} exhibit
substantial joint contribution, reflecting the well-documented interaction between
aggregator reliance and commission-driven margin erosion.

\subsection{Scenario Simulation: What-If Analysis}

The scenario simulation module permits real-time manipulation of all 13 input features
through interactive sliders. The quantile regression interval model provides a 95\%
prediction band, and a scenario risk classifier labels outcomes as \textit{Low},
\textit{Medium}, or \textit{High} risk based on interval width and presence of a
zero-crossing:

\begin{itemize}[leftmargin=*, itemsep=1pt]
  \item \textbf{Low Risk}: interval width $< 40\%$ of predicted value and lower bound $> 0$.
  \item \textbf{Medium Risk}: interval width 40--80\% of predicted value.
  \item \textbf{High Risk}: interval width $\geq 80\%$ or lower bound crosses zero.
\end{itemize}

Illustrative scenarios demonstrate that sustained profitability requires either
maintaining strong in-store and self-delivery shares or negotiating aggregator
commission rates well below the 30\% threshold.

\subsection{Prescriptive Optimisation Results}

SLSQP convergence was achieved in all tested profiles within 200 iterations. The
optimisation consistently recommends a higher in-store share and reduced reliance on
aggregator channels when COGS and OPEX rates are elevated, confirming the prescriptive
validity of the model. Average optimisation uplift ranges from 5\% to 15\% depending
on current channel configuration and restaurant segment.

% ────────────────────────────────────────────────────────────
\section{Strategic Recommendations}
\label{sec:recommendations}

Based on the predictive and optimisation outputs, the following evidence-based
strategic recommendations are offered to SkyCity Auckland management:

\begin{enumerate}[leftmargin=*, label=\textbf{R\arabic*.}, itemsep=2pt]

  \item \textbf{Prioritise In-Store and Self-Delivery Channels.}
  In-store dining consistently delivers the highest average net profit per month
  (\$2{,}257.32). Strategies to grow in-store traffic—loyalty programmes, dining
  events, and in-precinct promotions—should be prioritised. Complementary investment
  in self-delivery infrastructure is warranted, given 94.7\% of outlets are profitable
  within established cost guardrails.

  \item \textbf{Commission Rate Negotiation as Strategic Priority.}
  The empirically derived break-even commission rate is approximately 30\%. Any
  renegotiation with Uber Eats or DoorDash should target a rate at or below this
  threshold. A 1-percentage-point reduction in commission rate translates to an
  estimated \$200--\$300 monthly improvement in net profit for an average outlet.

  \item \textbf{Self-Delivery Expansion with Cost Guardrails.}
  Self-delivery operations remain profitable up to \$5.31 per order cost and 18 km
  delivery radius. 94.7\% of the current portfolio already operates within these bounds.
  Controlled expansion and operational improvements targeting cost reduction can unlock
  incremental margin recovery.

  \item \textbf{Segment and Subregion Resource Allocation.}
  Ghost Kitchen and Pizza concepts represent the highest-return strategic clusters, with
  net profits of \$9{,}051 and \$7{,}305 respectively. West Auckland similarly outperforms
  other regions. Capital allocation, marketing budgets, and new-opening decisions should be
  weighted toward these high-performance groupings.

  \item \textbf{Cost Structure Discipline.}
  Restaurants maintaining combined COGS + OPEX below 70\% of revenue consistently
  outperform portfolio average. Management should target COGS $\leq 28\%$ and OPEX
  $\leq 40\%$ as operational benchmarks.

  \item \textbf{Adopt Predictive Model for Ongoing Forecasting.}
  The XGBoost model (R\textsuperscript{2} = 0.9774) should replace static monthly reporting
  as the primary forecasting tool, with quarterly retraining on updated data to capture
  structural changes in demand and cost parameters.

\end{enumerate}

% ────────────────────────────────────────────────────────────
\section{Dashboard and Deployment}
\label{sec:dashboard}

The analytical pipeline was operationalised as a multi-page Streamlit web application
comprising five navigation modules:

\begin{enumerate}[leftmargin=*, itemsep=1pt]
  \item \textbf{Overview}: Portfolio KPIs, revenue-by-cuisine, profit-by-segment, channel
        mix composition charts, and subregional performance benchmarks.
  \item \textbf{Exploratory Analysis}: Distribution plots, correlation matrices, outlier
        identification, segment and subregion stratified analysis, and comprehensive
        trend visualisation.
  \item \textbf{Predictive Models}: Model training on demand, comparative performance
        tables, feature importance charts, predicted-vs-actual scatter plots, and
        cross-validation diagnostics.
  \item \textbf{What-If Simulator}: Interactive sliders for all 13 features, multi-model
        predictions with 95\% prediction intervals, risk classification, PDF scenario
        report generation, and multi-scenario comparison capability.
  \item \textbf{Optimisation Panel}: Commission sensitivity analysis, self-delivery
        break-even heatmaps, SLSQP channel-mix optimisation, and segment-stratified
        optimisation scoring.
\end{enumerate}

The application is deployable via \texttt{streamlit run app.py} and automates model
retraining when requested. All trained model objects are persisted to disk via
\texttt{joblib} and loaded on subsequent sessions to eliminate retraining latency
during interactive use.

% ────────────────────────────────────────────────────────────
\section{Limitations and Future Research}
\label{sec:limitations}

Several limitations should be acknowledged. First, the dataset covers a single
integrated precinct (SkyCity Auckland); generalisation to independent restaurants or
other geographic markets requires external validation. Second, temporal dynamics are
not modelled: the present framework treats each record as an independent observation
and does not capture seasonal patterns, promotional effects, or auto-regressive profit
dynamics. Third, the label-encoding approach for categorical variables assumes an
ordinal relationship among categories that may not hold; future work should evaluate
one-hot encoding or embedding layers.

Future research directions include: (\textit{i}) longitudinal modelling of profit
dynamics using recurrent architectures (LSTM, temporal Gradient Boosting);
(\textit{ii}) causal inference approaches to isolate the effect of commission rate
changes from confounding demand shifts; (\textit{iii}) multi-objective optimisation
incorporating fairness constraints across subregions; and (\textit{iv}) integration
with real-time point-of-sale feeds to enable live adaptive forecasting and automated
retraining pipelines.

% ────────────────────────────────────────────────────────────
\section{Conclusion}
\label{sec:conclusion}

This paper has presented a complete predictive and prescriptive intelligence framework
for multi-channel restaurant profit optimisation, applied to SkyCity Auckland
Restaurants \& Bars. Four machine-learning models were developed and rigorously
evaluated, with the best ensemble model (XGBoost) achieving an R\textsuperscript{2}
of 0.9774 and RMSE of NZD 852.71 on unseen data. The framework introduces quantile-based
prediction intervals for risk-stratified scenario planning and a SLSQP-based channel-mix
optimiser that identifies profit-maximising allocations within empirically grounded
constraints.

Key substantive findings include: (\textit{i}) in-store dining remains the most
profitable channel across all scenarios, generating \$2{,}257.32 average monthly net
profit; (\textit{ii}) a 30\% commission rate constitutes the economic break-even
threshold for third-party aggregator engagement; (\textit{iii}) self-delivery is viable
within defined cost (\$5.31 per order) and radius (18 km) guardrails, with 94.7\% of
the portfolio currently profitable; and (\textit{iv}) the majority of profit variance
(R\textsuperscript{2} = 0.9774) is explained by a 13-feature predictive model,
demonstrating that data-driven forecasting is both feasible and highly valuable in
restaurant operations.

By embedding these findings in an interactive decision-support dashboard, the study
demonstrates a pathway from static financial reporting to real-time, optimisation-driven
strategic management—a transition with significant potential implications for the
broader hospitality industry and multi-channel food service operators globally.

% ────────────────────────────────────────────────────────────
\section*{Acknowledgements}
The authors acknowledge the SkyCity Auckland management team for granting access to
operational data and providing domain expertise guidance throughout the analysis. We
thank the analytics and finance teams for data validation and contextual support.

% ────────────────────────────────────────────────────────────
\section*{Declaration of Competing Interests}
The authors declare no conflicts of interest.

% ────────────────────────────────────────────────────────────
\section*{Author Contributions}
Research Analytics Team: Conceptualisation, Methodology, Software Development, Data
Curation, Formal Analysis, Validation, Writing -- Original Draft, Writing -- Review
\& Editing, Visualisation.

% ────────────────────────────────────────────────────────────
% REFERENCES
% ────────────────────────────────────────────────────────────
\bibliographystyle{apalike}

\begin{thebibliography}{99}

\bibitem[Bertsimas and Kallus(2020)]{bertsimas2019prescrip}
Bertsimas, D., \& Kallus, N. (2020).
From predictive to prescriptive analytics.
\textit{Management Science}, 66(3), 1025--1044.

\bibitem[Breidbach et al.(2020)]{breidbach2020platform}
Breidbach, C.~F., Choi, S., Ellway, B., Keating, B.~W., Kettinger, W.,
Kuk, G., \ldots, \& Tumbas, S. (2020).
Operating without operations: How is technology changing the role of the firm?
\textit{Journal of Service Management}, 31(5), 1008--1033.

\bibitem[Cachon et al.(2022)]{cachon2022ownership}
Cachon, G.~P., Daniels, K.~M., \& Lobel, R. (2022).
The role of surge pricing on a service platform with self-scheduling capacity.
\textit{Manufacturing \& Service Operations Management}, 19(3), 368--384.

\bibitem[Chen and Guestrin(2016)]{chen2016xgboost}
Chen, T., \& Guestrin, C. (2016).
XGBoost: A scalable tree boosting system.
In \textit{Proceedings of the 22nd ACM SIGKDD International Conference on
Knowledge Discovery and Data Mining} (pp. 785--794). ACM.

\bibitem[Kwong et al.(2020)]{kwong2020menu}
Kwong, R., Kwong, F.~Y., \& Lam, M.~C. (2020).
Menu engineering and profitability analysis in Chinese restaurants.
\textit{International Journal of Contemporary Hospitality Management}, 32(7), 2487--2508.

\bibitem[Okumus et al.(2020)]{okumus2020revenue}
Okumus, B., Koseoglu, M.~A., \& Ma, F. (2020).
Food and gastronomy research in tourism and hospitality: A bibliometric analysis.
\textit{International Journal of Hospitality Management}, 73, 64--74.

\bibitem[Talluri and van Ryzin(2004)]{talluri2004theory}
Talluri, K.~T., \& Van Ryzin, G.~J. (2004).
\textit{The Theory and Practice of Revenue Management}.
Springer.

\bibitem[Wang et al.(2022)]{wang2022lastmile}
Wang, X., Zhan, L., Ruan, J., \& Zhang, J. (2022).
How to choose ``last mile'' delivery modes for e-commerce distribution:
A case study from China.
\textit{Sustainability}, 6(11), 6928.

\bibitem[Zhong et al.(2020)]{zhong2020forecasting}
Zhong, L., Guo, Y., Morrison, A.~M., \& Coca-Stefaniak, J.~A. (2020).
Links between tourists' urban experiences and place attachment:
Evidence from two heritage cities.
\textit{International Journal of Tourism Research}, 23(4), 608--623.

\end{thebibliography}

\end{document}
